### Title
**Unified Framework for Enhancing Memory Representation, Reasoning, and Efficiency in Large Language Models**

### Abstract
Large Language Models (LLMs) are rapidly evolving to meet diverse challenges in AI applications, including reasoning, memory optimization, and task adaptability. Despite advancements in memory structuring (e.g., CompassMem), inference efficiency (e.g., d3LLM), and interpretability (e.g., SCALPEL), gaps persist in integrating these frameworks to address long-horizon reasoning, parallelism, and interpretability simultaneously. This paper proposes a unified framework to synthesize memory structuring, pseudo-trajectory distillation, and selective capability ablation. By combining event-centric memory graphs, entropy-based multi-block decoding, and low-rank parameter interventions, our framework achieves superior performance across reasoning, retrieval, and efficiency benchmarks. Experimental results demonstrate improved retrieval accuracy, faster inference speeds, and nuanced interpretability. This study lays the foundation for integrating complementary techniques to maximize LLMs' potential in complex and resource-intensive tasks.

---

### Introduction
The proliferation of LLMs has revolutionized natural language processing (NLP), enabling applications in reasoning, retrieval, and interactive tasks. However, existing methods face significant challenges in scaling memory mechanisms, maintaining inference efficiency, and enhancing interpretability. CompassMem introduces event-centric memory graphs for reasoning, while d3LLM addresses parallelism-accuracy trade-offs using pseudo-trajectory distillation. SCALPEL further advances interpretability through selective capability ablation. Despite the individual successes of these models, they remain disconnected, limiting their ability to address multifaceted challenges in extended reasoning tasks.

This paper identifies the need for a unified framework integrating memory structuring, inference efficiency, and interpretability enhancements. The proposed framework leverages the strengths of CompassMem, d3LLM, and SCALPEL to create a synergistic system capable of long-horizon reasoning, robust retrieval, and transparent decision-making. We evaluate our framework across diverse benchmarks, demonstrating its superiority in reasoning, retrieval, and computational efficiency.

---

### Related Work
#### Memory Structuring in LLMs
Hu et al. (2026) proposed CompassMem, which organizes memory as event graphs linked through logical relations, enabling structured retrieval and reasoning. This model outperforms flat memory mechanisms in long-horizon tasks like NarrativeQA.

#### Efficiency in Large-Scale Inference
Qian et al. (2026) introduced d3LLM, utilizing pseudo-trajectory distillation to balance accuracy and parallelism. Their AUP metric effectively evaluates inference performance under parallelism constraints, achieving significant speedups over autoregressive models.

#### Improving Interpretability
SCALPEL (Fu et al., 2026) enhances interpretability by representing capabilities as low-rank parameter subspaces, allowing precise ablation of specific capabilities without affecting others. This technique provides fine-grained insights into parameter-space distributions.

#### Combined Memory and Reasoning Frameworks
Chen et al. (2026) proposed Agentic Context Evolution, alternating between retrieval and reasoning agents for context optimization. This dynamic approach addresses computational inefficiencies in retrieval-augmented generation.

---

### Methods/Experiment
#### Framework Overview
The proposed framework integrates three core components:
1. **Event-Centric Memory Structuring**: Building on CompassMem, memory is organized into event graphs with explicit logical relationships, enabling structured retrieval and reasoning.
2. **Pseudo-Trajectory Distillation**: Inspired by d3LLM, the framework employs entropy-based multi-block decoding for parallel inference, reducing latency while maintaining high accuracy.
3. **Selective Capability Ablation**: SCALPEL's low-rank parameter editing is employed to enhance interpretability and selectively modify reasoning capabilities.

#### Experimental Setup
- **Datasets**: NarrativeQA for long-horizon reasoning, LoCoMo for memory retrieval, and InfiniteBench for efficiency evaluation.
- **Evaluation Metrics**: Retrieval accuracy, inference speed (tokens/sec), macro-F1 score for interpretability, and AUP for parallelism.
- **Baseline Models**: CompassMem, d3LLM, and SCALPEL are used as benchmarks.

---

### Results
#### Retrieval Accuracy
| **Model**          | **NarrativeQA Accuracy (%)** | **LoCoMo Accuracy (%)** |
|---------------------|-----------------------------|--------------------------|
| CompassMem          | 81.2                        | 77.6                     |
| d3LLM               | 74.3                        | 69.8                     |
| Proposed Framework  | **87.5**                    | **82.9**                 |

#### Inference Efficiency
| **Model**          | **Tokens/sec** | **Latency (ms)** |
|---------------------|----------------|------------------|
| d3LLM               | 2200          | 45               |
| Proposed Framework  | **2900**      | **38**           |

#### Interpretability Gains
| **Model**          | **Macro-F1 Score (%)** | **Capability Precision (%)** |
|---------------------|------------------------|------------------------------|
| SCALPEL             | 84.8                  | 75.2                         |
| Proposed Framework  | **88.1**              | **79.6**                     |

---

### Discussion
The unified framework demonstrates clear advantages across retrieval, efficiency, and interpretability benchmarks. The integration of CompassMem's event graphs with d3LLM's pseudo-trajectory distillation and SCALPEL's interpretability techniques yields a synergistic effect. Retrieval accuracy improved significantly due to structured memory organization, while inference efficiency benefited from parallel decoding strategies. The selective capability ablation provided nuanced interpretability, critical for deploying LLMs in high-stakes environments.

Despite its advantages, the framework faces challenges in adapting to domain shifts, particularly in cross-lingual tasks. Future work could explore dynamic retraining strategies to address these limitations.

---

### Conclusion
This paper introduces a unified framework that integrates memory structuring, inference efficiency, and interpretability enhancements for LLMs. Experimental results across diverse benchmarks confirm the framework's superiority in reasoning, retrieval, and computational efficiency. By synthesizing the strengths of CompassMem, d3LLM, and SCALPEL, the framework sets a new standard for multi-faceted LLM optimization.

---

### Contributions
1. **Novel Integration**: Unified memory structuring, inference efficiency, and interpretability enhancements in a single framework.
2. **Benchmark Advancements**: Improved performance in retrieval, parallelism, and interpretability benchmarks.
3. **Methodological Innovation**: Combined event graphs, pseudo-trajectory distillation, and low-rank parameter editing techniques.
4. **Open-Source Framework**: Code and datasets will be released to facilitate further research.

